\documentclass[conference]{IEEEtran}
\IEEEoverridecommandlockouts
\usepackage{cite}
\usepackage{amsmath,amssymb,amsfonts}
\usepackage{algorithmic}
\usepackage{graphicx}
\usepackage{textcomp}
\usepackage{xcolor}
\usepackage[utf8]{inputenc}
\usepackage[spanish,es-nodecimaldot,es-noshorthands]{babel}
\usepackage{booktabs}
\usepackage{hyperref}

\def\BibTeX{{\rm B\kern-.05em{\sc i\kern-.025em b}\kern-.08em
    T\kern-.1667em\lower.7ex\hbox{E}\kern-.125emX}}

\begin{document}

\title{Sistema de Gestión y Seguimiento de Tareas para Mitigar Retrasos en Trabajos Colaborativos}

\author{
  \begin{tabular}{c@{\hspace{0.9cm}}c@{\hspace{0.9cm}}c}
    \begin{minipage}[t]{5.3cm}\centering
      \textbf{Denny Joaquin Mamani Huayhua}\\[0.08cm]
      \textit{Escuela Profesional de Ingeniería de Sistemas}\\
      \textit{Universidad Nacional de San Agustín de Arequipa}\\
      Arequipa, Perú\\
      \texttt{\small dmamanihuayh@unsa.edu.pe}
    \end{minipage} &
    \begin{minipage}[t]{5.3cm}\centering
      \textbf{Eder Lucio Vicente Medina}\\[0.08cm]
      \textit{Escuela Profesional de Ingeniería de Sistemas}\\
      \textit{Universidad Nacional de San Agustín de Arequipa}\\
      Arequipa, Perú\\
      \texttt{\small evicentem@unsa.edu.pe}
    \end{minipage} &
    \begin{minipage}[t]{5.3cm}\centering
      \textbf{Marco Antonio Salas Zegarra}\\[0.08cm]
      \textit{Escuela Profesional de Ingeniería de Sistemas}\\
      \textit{Universidad Nacional de San Agustín de Arequipa}\\
      Arequipa, Perú\\
      \texttt{\small msalasze@unsa.edu.pe}
    \end{minipage} \\[0.75cm]
    \begin{minipage}[t]{5.3cm}\centering
      \textbf{Rodrigo Enrique Neyra Chambilla}\\[0.08cm]
      \textit{Escuela Profesional de Ingeniería de Sistemas}\\
      \textit{Universidad Nacional de San Agustín de Arequipa}\\
      Arequipa, Perú\\
      \texttt{\small rneyra@unsa.edu.pe}
    \end{minipage} &
    \begin{minipage}[t]{5.3cm}\centering
      \textbf{Ronald David Jala Paricanaza}\\[0.08cm]
      \textit{Escuela Profesional de Ingeniería de Sistemas}\\
      \textit{Universidad Nacional de San Agustín de Arequipa}\\
      Arequipa, Perú\\
      \texttt{\small rjala@unsa.edu.pe}
    \end{minipage} &
    \begin{minipage}[t]{5.3cm}\centering
      \textbf{Mg. William Milton Bornas Rios}\\[0.08cm]
      \textit{Docente Asesor - EPIS}\\
      \textit{Universidad Nacional de San Agustín de Arequipa}\\
      Arequipa, Perú\\
      \texttt{\small wbornas@unsa.edu.pe}
    \end{minipage}
  \end{tabular}
}

\maketitle

\begin{abstract}
El desarrollo de trabajos grupales en entornos académicos y formativos enfrenta serias dificultades operativas cuando se emplean canales de mensajería instantánea como WhatsApp, evidenciándose saturación de mensajes en un 100\% de los casos, confusión de roles en un 76\% y retrasos en entregas en más del 64\%. Este artículo propone un sistema web fullstack para la gestión, seguimiento y mitigación de retrasos en proyectos colaborativos. La solución se fundamenta en el paradigma orientado a objetos y en estructuras de datos, articulando un entorno dual que integra Tableros Kanban para la gestión del flujo operativo y Diagramas de Gantt para la planificación temporal y detección temprana de cuellos de botella y alertas de retraso. La metodología sigue un ciclo de prototipado incremental en dos etapas: análisis y modelado funcional, seguido de la implementación modular y evaluación del prototipo. El sistema integra una arquitectura web desacoplada con backend modular en FastAPI y frontend reactivo y ligero. Los resultados preliminares y el diseño estructurado evidencian una reducción significativa en la incertidumbre de roles, alta claridad en la priorización de tareas y una gestión oportuna de los plazos de entrega en beneficio de equipos multidisciplinarios y la comunidad educativa.
\end{abstract}

\begin{IEEEkeywords}
aprendizaje colaborativo, estructuras de datos, tableros Kanban, diagrama de Gantt, prototipo web, seguimiento de tareas.
\end{IEEEkeywords}

\section{Introducción}
\subsection{Contexto y problemática}
En el ámbito de la educación superior, técnica y en equipos de investigación y desarrollo, el Aprendizaje Basado en Proyectos (PBL) y el trabajo colaborativo constituyen metodologías formativas esenciales. No obstante, investigaciones empíricas demuestran que el rendimiento de los equipos se fractura cuando no existe una estructuración formal de las actividades ni visibilidad sobre las prioridades y compromisos de trabajo entre pares, originando asimetrías de compromiso, conflictos de prioridad y retrasos no mitigados \cite{hussein2021addressing}. Para coordinar sus proyectos, los equipos recurren a plataformas de comunicación digital; sin embargo, al evaluar la experiencia de los usuarios, se comprueba que los entornos institucionales formales (como Microsoft Teams) suelen percibirse complejos o poco adaptables para proyectos cortos, propiciando que los grupos migren masivamente hacia canales de mensajería social desestructurados \cite{sobaih2021responses}.

En este contexto, la coordinación operativa de los equipos recae predominantemente en WhatsApp. Estudios empíricos en educación superior revelan que la informalidad comunicativa de dicha aplicación propicia la pérdida de información crítica en equipos numerosos ($\mu=4.12/5.00$), genera frecuentes malentendidos y conflictos por instrucciones ambiguas ($\mu=4.16/5.00$), y dificulta alcanzar consensos organizativos virtuales ($\mu=4.10/5.00$) \cite{vilches2019limitaciones}. Como consecuencia directa de esta desorganización, el 100\% de los estudiantes universitarios reconoce experimentar sobrecarga y saturación por mensajes desestructurados, un 76\% manifiesta confusión recurrente respecto a la asignación de roles y un 64\% reporta haber omitido fechas límite de entrega en trabajos grupales \cite{akram2024gams}.

Esta problemática trasciende el aula y se manifiesta en semilleros de investigación, proyectos multidisciplinarios, iniciativas de innovación y equipos de trabajo colaborativo que cursan múltiples entregas simultáneas. La falta de trazabilidad ocasiona entregas precipitadas de último minuto, desbalance en la carga laboral del equipo, fricciones interpersonales y estrés. En este escenario, el problema radica en la ausencia de una herramienta formal, intuitiva y accesible que gestione el flujo secuencial de actividades, visualice los plazos y estados de avance, y alerte oportunamente sobre retrasos y cuellos de botella en los proyectos colaborativos.

\subsection{Justificación}
La coordinación del trabajo colaborativo se debilita cuando carece de una estructura formal y de visibilidad sobre las prioridades y los compromisos de cada integrante, lo que genera ambigüedad de roles, conflictos de prioridad y retrasos no mitigados \cite{hussein2021addressing}. Cuando además la coordinación se apoya en canales de mensajería informal, la información crítica se dispersa y se pierde, y las instrucciones ambiguas producen malentendidos \cite{vilches2019limitaciones}. Aunque la evidencia disponible proviene del ámbito universitario, la misma dinámica es previsible en cualquier grupo pequeño que comparte un objetivo y coordina sus actividades por canales no estructurados, como equipos de voluntariado, de organización de eventos, de emprendimiento o de trabajo profesional. Esta extensión se plantea como hipótesis de diseño y solo se contrasta parcialmente en el piloto de este trabajo.

El proyecto es necesario porque ofrece un espacio centralizado y simple donde cada tarea tiene responsable, plazo y dependencias, y donde el sistema advierte oportunamente cuando un retraso puede afectar al resto del trabajo. Con ello sustituye la informalidad comunicativa por trazabilidad, claridad de responsabilidades y control del avance.

\textbf{Responsabilidad social.} La herramienta promueve la transparencia y la equidad en el reparto de la carga de trabajo, reduce el desgaste relacional y el estrés que provoca la descoordinación, y fortalece competencias de responsabilidad compartida. Su uso no depende de conocimientos técnicos ni de costos, por lo que puede beneficiar a grupos con recursos limitados.

\textbf{Innovación.} La propuesta integra el tablero Kanban con la precedencia entre tareas: una tarea cuyas predecesoras no han terminado se marca como bloqueada, y las alertas de retraso se calculan según su impacto sobre las tareas que dependen de ella. Todo ello en una aplicación web de arranque guiado, adaptable a distintos tamaños de pantalla, incluido el celular.

\subsection{Objetivos}
\textbf{Objetivo general.} Desarrollar y evaluar, hasta el 7 de diciembre de 2026, un prototipo de aplicación web de uso general, adaptable a pantallas de celular y de escritorio, para la gestión y el seguimiento del trabajo colaborativo, que combine un tablero Kanban con dependencias entre tareas y alertas de retraso, a fin de mitigar retrasos y confusión de responsabilidades en grupos de personas.

\textbf{Objetivos específicos.}
\begin{itemize}
    \item \textbf{OE1.} Analizar los requisitos funcionales del trabajo colaborativo y elaborar su modelo funcional (casos de uso, flujo de actividades, modelo de datos y reglas de bloqueo y alertas) durante la Etapa~I.
    \item \textbf{OE2.} Diseñar la arquitectura del sistema (React y Supabase) y el módulo de lógica de dominio que calcula el estado de las tareas y las alertas durante la Etapa~I.
    \item \textbf{OE3.} Implementar, durante la Etapa~II, un prototipo funcional con el alcance delimitado en la Sección~\ref{sec:alcance}.
    \item \textbf{OE4.} Evaluar, al cierre de la Etapa~II, el funcionamiento y la usabilidad del prototipo mediante casos de prueba y la escala System Usability Scale (SUS) con al menos 10 usuarios, y contrastar el puntaje obtenido con el valor de referencia de 68 puntos.
\end{itemize}

\subsection{Alcance y limitaciones}
\label{sec:alcance}
\textbf{Delimitación.} El prototipo incluye:
\begin{itemize}
    \item acceso con cuenta de Google, creación de proyectos y unión mediante código o enlace de invitación;
    \item tareas con responsable, plazo, estado y precedencia entre tareas;
    \item tablero Kanban usable en pantalla de celular, con bloqueo automático de tareas según sus predecesoras;
    \item alertas de retraso calculadas al ingresar y mostradas dentro de la aplicación;
    \item notas de estado al solicitar apoyo o modificar un plazo;
    \item un cronograma resumido de solo lectura para las tareas con precedencia, en caso de que el tiempo lo permita.
\end{itemize}

Quedan fuera del alcance: mensajería o comentarios libres, archivos adjuntos, notificaciones externas o programadas, videoconferencia, integración con plataformas de mensajería, aplicación móvil nativa y funciones específicas de un dominio profesional particular.

\textbf{Plan del proyecto.} El proyecto se organiza en dos etapas, cuyas actividades y plazos se resumen en el Cuadro~\ref{tab:plan_proyecto}.

\begin{table}[htbp]
\caption{Plan de Actividades y Plazos del Proyecto}
\label{tab:plan_proyecto}
\centering
\begin{tabular}{cp{4.4cm}p{2.3cm}}
\toprule
\textbf{Etapa} & \textbf{Actividad} & \textbf{Plazo} \\
\midrule
I & Diagnóstico del problema y revisión de literatura & Octubre 2026 \\
I & Definición de la propuesta y retroalimentación del docente & Octubre--noviembre 2026 \\
I & Replanteamiento hacia un uso general & Noviembre 2026 \\
I & Requisitos, modelo funcional y arquitectura & 25--26 de noviembre \\
II & Implementación del prototipo & 27 de noviembre--3 de diciembre \\
II & Pruebas y evaluación con usuarios (SUS) & 4--5 de diciembre \\
II & Análisis de resultados, redacción final y entrega & 5--7 de diciembre \\
\bottomrule
\end{tabular}
\end{table}

\textbf{Recursos.} El equipo está conformado por cinco integrantes y un docente asesor, con la siguiente distribución de responsabilidades: desarrollo del prototipo (interfaz, integración con Supabase y lógica de dominio), modelado funcional y de datos, diseño de la interfaz, pruebas y evaluación de usabilidad, y documentación y redacción del artículo. Se emplean Git y GitHub para el control de versiones, React con TypeScript para la interfaz, Supabase como backend y base de datos, Draw.io para el modelado, y agentes de inteligencia artificial bajo un proceso de desarrollo guiado por especificaciones.

\textbf{Limitaciones.} La validación se realiza con una muestra piloto pequeña y en un periodo corto, por lo que los resultados son exploratorios y no permiten afirmar un efecto estadístico sobre la reducción de retrasos. La aplicabilidad a dominios distintos del educativo se plantea como hipótesis de diseño y no se evalúa de forma separada. El prototipo prioriza el funcionamiento esencial y deja para trabajos futuros las funciones excluidas en la delimitación.

El artículo se organiza de la siguiente manera: la Sección II expone el estado del arte y el análisis comparativo de la literatura; la Sección III describe la metodología adoptada; la Sección IV detalla el análisis y modelamiento funcional con estructuras de datos; la Sección V presenta el desarrollo del prototipo; la Sección VI resume las pruebas y resultados; la Sección VII discute los hallazgos; y la Sección VIII presenta las conclusiones y líneas de trabajo futuro.

\section{Estado del Arte}
\subsection{Marco conceptual}
El trabajo colaborativo mediado por tecnologías educativas requiere estructuración de roles y sincronización de entregables para evitar fallas organizativas \cite{hussein2021addressing, sobaih2021responses}. Para mitigar el desorden derivado de canales informales \cite{vilches2019limitaciones}, la literatura ha explorado Sistemas de Información Gerencial (MIS) adaptados a la docencia que centralizan cronogramas y roles \cite{akram2024gams}. Asimismo, se han desarrollado herramientas de administración de actividades enfocadas en la priorización de plazos para estudiantes individuales \cite{nurzi2022task}, así como Sistemas de Soporte a Decisiones Grupales (GDSS) orientados a automatizar la planificación académica y reducir el error humano frente a registros manuales \cite{siddiqui2018gdss}.

Para medir la aceptación y efectividad de estos sistemas de software, la literatura de interacción humano-computador adopta el System Usability Scale (SUS) como instrumento estandarizado y robusto de evaluación de usabilidad \cite{brooke1996sus}. Por otra parte, desde la perspectiva algorítmica y de estructuras de datos, la gestión ágil fundamenta el uso de Tableros Kanban para limitar el trabajo en progreso ($WIP$) y optimizar el flujo de tareas, combinados con Diagramas de Gantt para la representación temporal de plazos y fechas límite. Sobre estas estructuras, el monitoreo dinámico permite detectar cuellos de botella en columnas críticas y emitir alertas preventivas cuando una actividad supera su fecha estimada de entrega.

\subsection{Trabajos relacionados}
La revisión sistemática de la literatura académica reciente evidencia diversas aproximaciones para atender la gestión de proyectos en la educación superior:
\begin{itemize}
    \item Hussein \cite{hussein2021addressing} analizó los factores determinantes del fracaso en proyectos de ingeniería colaborativos mediante 67 reportes de reflexión estudiantil, identificando que la ambigüedad en la asignación de roles y los conflictos de prioridad son las causas raíz del bajo rendimiento, recomendando la estructuración formal del flujo de trabajo.
    \item Sobaih et al. \cite{sobaih2021responses} evaluaron la experiencia de aprendizaje comparando plataformas institucionales formales (Microsoft Teams) con redes sociales, demostrando que las herramientas corporativas presentan curvas de aprendizaje elevadas y redundancia funcional que desincentivan la adopción rápida en grupos universitarios.
    \item Vilches Vilela y Reche Urbano \cite{vilches2019limitaciones} condujeron un estudio empírico cuantitativo con 1709 universitarios, demostrando que el uso de WhatsApp como canal primario de coordinación propicia dispersión informativa, pérdida de archivos clave ($\mu=4.12/5.00$) y tensiones interpersonales por instrucciones ambiguas ($\mu=4.16/5.00$). Su valor radica en el diagnóstico cuantitativo, careciendo de una propuesta de software.
    \item Akram et al. \cite{akram2024gams} desarrollaron GAMS, un sistema web en PHP y MySQL con perfiles para administradores, líderes y miembros, incorporando recordatorios vía SMS y alcanzando un puntaje SUS de 70.44. Sin embargo, su arquitectura gestiona tareas como registros tabulares independientes, sin mecanismos visuales e interactivos de flujo ágil de trabajo.
    \item Nurzi y Ab Wahab \cite{nurzi2022task} presentaron un sistema web de seguimiento de tareas estudiantiles centrado en la priorización por fecha de vencimiento. Aunque mitiga olvidos individuales, el sistema fue concebido para uso monousuario, omitiendo la coordinación colaborativa y la interdependencia entre compañeros.
    \item Siddiqui et al. \cite{siddiqui2018gdss} implementaron un sistema de soporte a decisiones grupales para la planificación de periodos académicos, validando que los flujos automatizados en bases de datos eliminan errores de coordinación y reducen tiempos de gestión frente a hojas de cálculo manuales; no obstante, su alcance está orientado al nivel macro administrativo.
\end{itemize}

\subsection{Análisis comparativo}
En la Tabla~\ref{tab:comparativa} se sintetizan las características, aportes y limitaciones de las soluciones identificadas en la literatura frente a la presente propuesta.

\begin{table}[htbp]
\caption{Comparación de Soluciones en el Estado del Arte}
\label{tab:comparativa}
\centering
\begin{tabular}{p{1.4cm}p{1.7cm}p{2.3cm}p{1.7cm}}
\toprule
\textbf{Referencia} & \textbf{Enfoque / Tech} & \textbf{Aporte Principal} & \textbf{Limitación} \\
\midrule
Hussein \cite{hussein2021addressing} & Cualitativo (PBL) & Diagnóstico de causas de fricción en ingeniería & Enfoque formativo sin software \\
Sobaih et al. \cite{sobaih2021responses} & Entornos formales & Comparativa Teams vs. redes sociales & Complejidad y baja adaptabilidad \\
Vilches et al. \cite{vilches2019limitaciones} & Empírico (1709 alumnos) & Evidencia estadística de fallas en mensajería & Diagnóstico sin artefacto técnico \\
Akram et al. \cite{akram2024gams} & Web PHP / MySQL & Roles y notificaciones SMS & Tareas planas sin flujo visual \\
Nurzi et al. \cite{nurzi2022task} & Web / Base de datos & Priorización temporal de tareas & Solo monousuario (individual) \\
Siddiqui et al. \cite{siddiqui2018gdss} & Soporte a decisiones (DSS) & Automatización de calendarización académica & Orientado al nivel macro administrativo \\
\midrule
\textbf{Propuesta} & \textbf{Fullstack (FastAPI + Web)} & \textbf{Flujo de 15 actividades, vista dual Kanban + Gantt, roles y detección de retrasos} & \textbf{Evaluación inicial en muestra piloto} \\
\bottomrule
\end{tabular}
\end{table}

\subsection{Brecha identificada y aporte de la propuesta}
A partir del análisis comparativo se constata una marcada brecha en el estado del arte: las herramientas empleadas en la actualidad se polarizan entre canales de mensajería desestructurados que propician saturación comunicativa \cite{vilches2019limitaciones}, suites corporativas de alta complejidad que dificultan la adopción en proyectos breves \cite{sobaih2021responses}, o sistemas que gestionan actividades como registros planos y aislados \cite{akram2024gams, nurzi2022task}. Como demostró Hussein \cite{hussein2021addressing}, la desarticulación en las prioridades de entrega vulnera la efectividad de los equipos. Ninguna de las soluciones revisadas integra un flujo formal orientado al trabajo colaborativo (las 15 actividades de coordinación) con una arquitectura dual que vincule tableros Kanban para el flujo operativo diario con diagramas de Gantt interactivos capaces de alertar tempranamente cuellos de botella y desviaciones de cronograma. La presente propuesta cubre esta necesidad articulando un modelado funcional riguroso con estructuras de datos en una arquitectura fullstack ligera y escalable.

\section{Metodología}
\subsection{Enfoque metodológico}
Para el desarrollo de la investigación y la construcción del artefacto tecnológico se adopta un enfoque de \textbf{Prototipado Incremental y Evolutivo}, correspondiente a proyectos de ingeniería de sistemas en etapas tempranas de formación (Categoría I). Este marco metodológico permite abordar de manera iterativa el modelado de estructuras de datos y componentes visuales, facilitando ciclos continuos de refinamiento funcional a partir de la verificación de requisitos y la retroalimentación de casos de uso representativos del trabajo en equipo.

\subsection{Fases y actividades}
El proceso del proyecto se organiza en dos etapas de ingeniería de software.

\textbf{Etapa I: Análisis y modelado.}
\begin{itemize}
    \item Diagnóstico del problema de coordinación en el trabajo colaborativo y revisión de literatura.
    \item Definición de la propuesta, retroalimentación del docente y replanteamiento hacia un uso general.
    \item Elicitación y priorización de requisitos funcionales.
    \item Elaboración del modelo funcional: actores y casos de uso, flujo de actividades del usuario, modelo de datos y reglas de negocio.
    \item Diseño de la arquitectura y del módulo de lógica de dominio.
\end{itemize}
\textit{Entregables:} requisitos priorizados, modelo funcional y especificación del sistema.

\textbf{Etapa II: Construcción y evaluación.}
\begin{itemize}
    \item Implementación incremental del prototipo mediante desarrollo guiado por especificaciones (SDD), con asistencia de agentes de inteligencia artificial.
    \item Pruebas unitarias de la lógica de dominio y pruebas funcionales de los flujos principales.
    \item Evaluación de usabilidad con al menos 10 usuarios mediante la escala SUS.
    \item Análisis de resultados y redacción del artículo.
\end{itemize}
\textit{Entregables:} prototipo funcional, resultados de pruebas y de usabilidad, y artículo final.

\subsection{Herramientas de apoyo}
\begin{itemize}
    \item \textbf{Gestión de código y colaboración:} Git y GitHub para el control de versiones y la trazabilidad del desarrollo.
    \item \textbf{Modelado:} Draw.io para la arquitectura, los diagramas UML y los flujogramas.
    \item \textbf{Plataforma tecnológica del prototipo:} React con TypeScript para la interfaz web adaptable a celular; Supabase (PostgreSQL, autenticación con Google y reglas de acceso por fila) como backend; y un módulo de lógica de dominio independiente de la interfaz y de la base de datos, que calcula bloqueos, ruta crítica y alertas.
    \item \textbf{Proceso de desarrollo:} desarrollo guiado por especificaciones con asistencia de agentes de inteligencia artificial.
    \item \textbf{Evaluación:} escala SUS \cite{brooke1996sus} y casos de prueba funcionales.
\end{itemize}

\section{Análisis y Modelamiento Funcional}
\subsection{Requisitos}
\subsection{Modelo funcional}
El modelo funcional describe el ciclo de trabajo que la aplicación guía al usuario, estructurado en 15 actividades agrupadas en tres macrofases. Es una guía y no un proceso obligatorio: cada grupo puede omitir los pasos que no necesite.

\textbf{Fase 1: Planificación (actividades 1 a 8).} Inicio formal del trabajo (1), creación del proyecto e invitación de integrantes (2), registro de integrantes (3), definición de acuerdos y objetivos compartidos (4), división del trabajo en tareas (5), asignación de responsables (6), establecimiento de plazos (7) y declaración de dependencias entre tareas con su organización en el tablero (8).

\textbf{Fase 2: Ejecución y seguimiento (actividades 9 a 11).} Ejecución de las tareas individuales (9), actualización del estado y del avance en el tablero (10) y revisión continua del progreso, con alertas y la posibilidad de solicitar apoyo ante bloqueos (11).

\textbf{Fase 3: Control y cierre (actividades 12 a 15).} Evaluación de retrasos y cuellos de botella (12), reorganización de plazos y reasignación de tareas críticas (13), revisión final del trabajo realizado (14) y cierre del proyecto con un resumen de aportes por integrante, visible solo para el grupo (15).

\textbf{Reglas principales.}
\begin{itemize}
    \item Una tarea con predecesoras no terminadas se considera bloqueada.
    \item No se permiten dependencias que formen ciclos.
    \item Las tareas con precedencia determinan una ruta crítica, es decir, la secuencia de tareas cuyo retraso retrasa el trabajo completo.
    \item Las alertas se calculan al ingresar a la aplicación y cubren tareas próximas a vencer o vencidas, tareas retrasadas que bloquean a otras, acumulación excesiva en una columna y desbalance de carga entre integrantes.
\end{itemize}

El prototipo implementa por completo las actividades 2, 3 y 5 a 13, y ofrece soporte básico para las actividades 1, 4, 14 y 15 mediante un campo de objetivo del proyecto, el cierre del proyecto y un resumen de tareas completadas por integrante.
\subsection{Diseño de la solución}

\section{Desarrollo del Prototipo}
\subsection{Tecnologías y herramientas}
\subsection{Implementación}
\subsection{Repositorio y demostración}

\section{Pruebas y Resultados}
\subsection{Casos de prueba}
\subsection{Resultados obtenidos}

\section{Discusión}

\section{Conclusiones y Trabajo Futuro}

\section*{Agradecimiento}

\bibliographystyle{IEEEtran}
\bibliography{referencias}

\end{document}
